% Source: transfer_learning.tex
\begin{notechapter}
% ---- chapter-local definitions ----
\lstdefinestyle{python}{
    language=Python,
    basicstyle=\ttfamily\small,
    keywordstyle=\color{blue},
    commentstyle=\color{green!50!black},
    stringstyle=\color{red},
    showstringspaces=false,
    breaklines=true,
    frame=single,
    numbers=left,
    numberstyle=\tiny\color{gray},
    captionpos=b
}
% -----------------------------------
\chapter{Transfer Learning}\label{ch:transfer-learning}

\section{Introduction}
Transfer learning is a powerful technique in machine learning and
deep learning, where a model developed for one task is reused as
the starting point for a model on a second task. It is particularly
useful when you have limited data for the second task or when
training a model from scratch is computationally expensive.

\section{What is Transfer Learning?}
Transfer learning involves taking a pre-trained model (usually
trained on a large dataset) and adapting it to a new, often related
task. Instead of training a model from scratch, you leverage the
knowledge (e.g., feature representations) learned by the pre-trained
model. This approach is especially popular in deep learning, where
models like convolutional neural networks (CNNs) require large
amounts of data and computational resources to train.

\section{Why Use Transfer Learning?}
\begin{itemize}
    \item \textbf{Saves Time and Resources}: Training deep learning
          models from scratch can take days or weeks, even on powerful
          hardware. Transfer learning reduces this time significantly.
    \item \textbf{Works with Small Datasets}: When you have limited
          labeled data for your task, transfer learning allows you to
          leverage a pre-trained model's knowledge.
    \item \textbf{Improves Performance}: Pre-trained models have already
          learned useful features (e.g., edges, textures, shapes) from
          large datasets, which can generalize well to new tasks.
\end{itemize}

\section{How Does Transfer Learning Work?}
Transfer learning typically involves the following steps:

\subsection{Choose a Pre-trained Model}
Select a model that has been trained on a large, general-purpose
dataset (e.g., ImageNet\footnote{ImageNet is a large-scale dataset of over 14 million labeled images, categorized into more than 20,000 classes. It is widely used for training and benchmarking deep learning models, particularly in computer vision. The ImageNet Large Scale Visual Recognition Challenge (ILSVRC) has driven significant advancements in image classification, object detection, and other vision tasks. Many pre-trained models (e.g., VGG, ResNet, Inception) are trained on ImageNet and achieve state-of-the-art performance.} for image tasks). Popular pre-trained models
include:
\begin{itemize}
    \item \textbf{VGG}: A deep CNN architecture with 16-19 layers,
          known for its simplicity and strong performance on image
          classification tasks.
    \item \textbf{ResNet}: A residual network with "skip connections"
          that allow training of very deep networks (e.g., ResNet-50,
          ResNet-101). It addresses the vanishing gradient problem.
    \item \textbf{Inception}: A family of models (e.g., InceptionV3)
          that use multi-scale convolutions to capture features at
          different resolutions.
    \item \textbf{EfficientNet}: A scalable CNN architecture that
          achieves state-of-the-art performance with fewer parameters
          and computations.
    \item \textbf{BERT}: A transformer-based model for natural language
          processing tasks, pre-trained on large text corpora.
\end{itemize}

\subsection{Freeze Early Layers}
The early layers of a pre-trained model typically learn low-level
features (e.g., edges, textures) that are general and applicable to
many tasks. These layers are often "frozen" (their weights are not
updated during training) to retain their learned features.

\subsection{Replace and Train the Final Layers}
The final layers of the pre-trained model are replaced with new layers
tailored to the new task (e.g., a new classification head for a different
number of classes). Only these new layers (and sometimes a few of the
later layers of the pre-trained model) are trained on the new dataset.

\subsection{Fine-Tune (Optional)}
After training the new layers, you can optionally "fine-tune" the model
by unfreezing some of the earlier layers and training them with a very
low learning rate. This allows the model to adapt more specifically to
the new task.

\section{Feature Extraction with a Custom Classifier}
In some cases, instead of replacing the final layers of a pre-trained model
with new neural network layers, you can replace them with a different
classifier, such as a \textbf{Support Vector Machine (SVM)}, \textbf{Random Forest},
or other machine learning algorithms. This approach is often referred to as
\textbf{feature extraction with a custom classifier} or \textbf{hybrid transfer learning}.

\subsection{How It Works}
\begin{itemize}
    \item \textbf{Step 1: Feature Extraction}:
          Use a pre-trained model to extract features from your dataset.
          Remove the final classification layer(s) of the pre-trained model,
          as these are typically task-specific (e.g., softmax for classification).
          Pass your input data through the pre-trained model to obtain high-level
          feature representations (feature vectors).
    \item \textbf{Step 2: Train a Custom Classifier}:
          Use the extracted features as input to train a custom classifier,
          such as an SVM, Random Forest, or Gradient Boosting Machine (GBM).
          The custom classifier learns to map the feature vectors to the target labels.
\end{itemize}

\subsection{Dimensionality Reduction}
When working with high-dimensional feature vectors (e.g., from a pre-trained CNN),
it is often beneficial to apply \textbf{dimensionality reduction techniques} before training
the custom classifier. This step can:
\begin{itemize}
    \item \textbf{Reduce Computational Cost}: High-dimensional feature vectors can be
          computationally expensive to process, especially for algorithms like SVMs.
    \item \textbf{Prevent Overfitting}: Reducing the number of features can help prevent
          overfitting, particularly when the dataset is small.
    \item \textbf{Improve Performance}: By retaining only the most important features,
          dimensionality reduction can improve the performance of the classifier.
\end{itemize}

Common dimensionality reduction techniques include:
\begin{itemize}
    \item \textbf{Principal Component Analysis (PCA)}: A linear technique that projects
          the data onto a lower-dimensional space while preserving as much variance as possible.
    \item \textbf{t-SNE}: A non-linear technique often used for visualization in 2D or 3D.
    \item \textbf{UMAP}: A modern alternative to t-SNE that is faster and preserves more global structure.
\end{itemize}

\subsection{Example: Using PCA for Dimensionality Reduction}
Below is an example of applying PCA to reduce the dimensionality of extracted features
before training an SVM classifier:

\begin{lstlisting}[language=Python, style=python]
from sklearn.decomposition import PCA
from sklearn.svm import SVC
from sklearn.model_selection import train_test_split
from sklearn.metrics import accuracy_score

# Assume `features` is the high-dimensional feature matrix from the
# pre-trained model and `labels` is the corresponding target labels.

# Apply PCA to reduce dimensionality
pca = PCA(n_components=100)  # Reduce to 100 dimensions
features_reduced = pca.fit_transform(features)

# Split the data into training and testing sets
X_train, X_test, y_train, y_test = train_test_split(features_reduced,
                                                    labels, test_size=0.2,
                                                    random_state=42)

# Train an SVM classifier
svm_classifier = SVC(kernel='linear')
svm_classifier.fit(X_train, y_train)

# Evaluate the classifier
y_pred = svm_classifier.predict(X_test)
accuracy = accuracy_score(y_test, y_pred)
print(f"Test Accuracy: {accuracy * 100:.2f}%")
\end{lstlisting}

\subsection{Advantages}
\begin{itemize}
    \item \textbf{Flexibility}: You can use any machine learning algorithm as the final classifier.
    \item \textbf{Interpretability}: Algorithms like decision trees or SVMs are often more interpretable than deep neural networks.
    \item \textbf{Efficiency}: Training a simpler classifier on extracted features is often faster and requires fewer computational resources.
    \item \textbf{Small Datasets}: This approach works well when you have a small dataset, as the pre-trained model provides robust feature representations.
\end{itemize}

\subsection{Example: Replacing the Last Layer with an SVM}
Below is an example of using a pre-trained ResNet50 model for feature extraction
and training an SVM classifier:

\begin{lstlisting}[language=Python, style=python]
import numpy as np
from tensorflow.keras.applications import ResNet50
from tensorflow.keras.layers import GlobalAveragePooling2D
from tensorflow.keras.models import Model
from sklearn.svm import SVC
from sklearn.model_selection import train_test_split
from sklearn.metrics import accuracy_score

# Load pre-trained ResNet50 model (without the top classification layer)
base_model = ResNet50(weights='imagenet', include_top=False,
                      input_shape=(224, 224, 3))

# Add a global average pooling layer to reduce the feature map to a vector
x = base_model.output
x = GlobalAveragePooling2D()(x)

# Create the feature extraction model
feature_extraction_model = Model(inputs=base_model.input, outputs=x)

# Load and preprocess your dataset
# Assume `images` is a numpy array of shape (num_samples, 224, 224, 3)
# and `labels` is a numpy array of shape (num_samples,)
images = np.random.rand(100, 224, 224, 3)  # Replace with your data
labels = np.random.randint(0, 2, 100)      # Replace with your labels

# Extract features using the pre-trained model
features = feature_extraction_model.predict(images)

# Split the data into training and testing sets
X_train, X_test, y_train, y_test = train_test_split(features, labels,
                                                    test_size=0.2,
                                                    random_state=42)

# Train an SVM classifier
svm_classifier = SVC(kernel='linear')
svm_classifier.fit(X_train, y_train)

# Evaluate the classifier
y_pred = svm_classifier.predict(X_test)
accuracy = accuracy_score(y_test, y_pred)
print(f"Test Accuracy: {accuracy * 100:.2f}%")
\end{lstlisting}

\subsection{Example: Replacing the Last Layer with a Random Forest}
Similarly, you can replace the last layer with a Random Forest classifier:

\begin{lstlisting}[language=Python, style=python]
from sklearn.ensemble import RandomForestClassifier

# Train a Random Forest classifier
rf_classifier = RandomForestClassifier(n_estimators=100, random_state=42)
rf_classifier.fit(X_train, y_train)

# Evaluate the classifier
y_pred = rf_classifier.predict(X_test)
accuracy = accuracy_score(y_test, y_pred)
print(f"Test Accuracy: {accuracy * 100:.2f}%")
\end{lstlisting}

\section{Applications of Transfer Learning}
Transfer learning is widely used in various domains:
\begin{itemize}
    \item \textbf{Computer Vision}:
          \begin{itemize}
              \item Image classification, object detection, segmentation.
              \item Example: Using a pre-trained ResNet model for medical
                    image diagnosis.
          \end{itemize}
    \item \textbf{Natural Language Processing (NLP)}:
          \begin{itemize}
              \item Text classification, sentiment analysis, machine translation.
              \item Example: Fine-tuning BERT for a specific text classification task.
          \end{itemize}
    \item \textbf{Audio Processing}:
          \begin{itemize}
              \item Speech recognition, sound classification.
          \end{itemize}
    \item \textbf{Reinforcement Learning}:
          \begin{itemize}
              \item Transferring policies from one environment to another.
          \end{itemize}
\end{itemize}

\section{Challenges and Considerations}
\begin{itemize}
    \item \textbf{Domain Mismatch}: If the new task is very different from
          the original task, the pre-trained model may not be effective.
    \item \textbf{Overfitting}: With a small dataset, the model may overfit
          to the training data. Techniques like data augmentation and
          regularization can help.
    \item \textbf{Choosing the Right Model}: Not all pre-trained models are
          suitable for every task. The choice depends on the complexity of
          the task and the available data.
\end{itemize}

\section{Conclusion}
Transfer learning leverages pre-trained models to save time, resources,
and improve performance on new tasks. It is particularly useful when you
have limited data or computational resources. Popular pre-trained models
include ResNet, VGG, Inception, and BERT. The process involves freezing
early layers, replacing final layers, and optionally fine-tuning. Additionally,
you can replace the final layers with custom classifiers like SVMs or Random
Forests for greater flexibility and interpretability. When working with
high-dimensional feature vectors, dimensionality reduction techniques like PCA
can further improve efficiency and performance.
\end{notechapter}
